% ============================================================
% Proposition 3.1: Local Existence and Uniqueness
% Status: CORRECTION (original C∞ claim is false)
% ============================================================

% --- Definitions and Assumptions ---
\begin{definition}[State Space]
Let $\mathcal{X} = [0, \Delta_{\max}]^{|D||M|} \times [0,1]^{3|D||M| + 3}$
with the Euclidean topology. Denote $\text{int}(\mathcal{X})$ as the
interior and $\partial\mathcal{X}$ as the boundary.
\end{definition}

\begin{definition}[Vector Field]
Define $\mathbf{f}: \mathcal{X} \to \R^N$ by collecting the right-hand
sides of the ODE system:
\begin{align*}
\dot{\delta}_A(d,m,t) &= f_{\text{learn}} \cdot \eta \cdot T_A
  - \lambda_c (1 - \gamma_\sigma \sigma_A) \delta_A (1 - r_A), \\
\dot{\sigma}_A(d,m,t) &=
  \rho P_A \alpha_A (1 - \gamma_\sigma \sigma_A) - \epsilon_\sigma \Omega_{SL} \sigma_A, \\
\dot{\alpha}_A(d,m,t) &=
  \gamma C_A (1 - \alpha_A) - \zeta_\alpha \alpha_A R_A^{\text{surface}}, \\
\dot{\hat{M}}_A(d,m,t) &=
  \nu_M (\sigma_A - \hat{M}_A) - \xi_M \Omega_{SL} \hat{M}_A, \\
\dot{\Xi}_A^P(t) &= \kappa_P (P^* - \Xi_A^P), \\
\dot{\Xi}_A^I(t) &= \kappa_I (I^* - \Xi_A^I), \\
\dot{\Xi}_A^F(t) &= \kappa_F (F^* - \Xi_A^F).
\end{align*}
The Stage~1 proxy $\tilde{\sigma}_A$ enters through the fusion equation:
\[
\tilde{\sigma}_A = w_g \max(0, g_A) + w_r \max(0, r_A) + w_c c_A,
\]
where $g_A$, $r_A$, $c_A$ involve $\Corr$ and $\CosSim$.
\end{definition}

\begin{assumption}[Regularity of Coefficient Functions]
\label{assum:regularity}
The coefficient functions satisfy:
\begin{enumerate}[nosep]
\item $P_A, C_A, \Omega_{SL}, R_A^{\text{surface}}, T_A, r_A, P^*, I^*, F^*$
  are uniformly bounded and continuous in $t$ for all $x \in \mathcal{X}$.
\item $P_A, C_A, \Omega_{SL}$ are locally Lipschitz in $x$ with common
  constant $L$.
\item $\Corr$ and $\CosSim$ are evaluated on finite sample batches;
  the resulting functions $\tilde{g}_A, \tilde{r}_A, \tilde{c}_A$ are
  piecewise-$C^\infty$ with finite jump discontinuities bounded by
  $O(1/\sqrt{n_{\text{batch}}})$ (see Proposition~3.6).
\end{enumerate}
\end{assumption}

% --- Proposition Statement ---
\begin{proposition}[Local Existence and Uniqueness]
\label{prop:local-existence}
Under Assumption~\ref{assum:regularity}, for any initial condition
$\mathbf{x}_A(0) \in \text{int}(\mathcal{X})$, the coupled ODE system
admits a unique solution $\mathbf{x}_A(t)$ defined on a maximal interval
$[0, T_{\max})$ where $T_{\max} > 0$.
\end{proposition}

% --- Proof ---
\begin{proof}
We proceed in three steps: (i) correct the smoothness claim of the
original proof, (ii) verify local Lipschitz continuity, and (iii)
apply Picard--Lindelöf.

\paragraph{Step 1: The vector field is not $C^\infty$.}
The original manuscript (line~853) claims all component functions are
$C^\infty$ on $\text{int}(\mathcal{X})$. This is false for two terms:
\begin{enumerate}[nosep]
\item The Stage~1 fusion (Eq.~11 in the manuscript, $\tilde{\sigma}_A = w_g \max(0, g_A) + w_r \max(0, r_A) + w_c c_A$) contains $\max(0, g_A)$,
  which is not differentiable at $g_A = 0$.
\item The RGA signal (Eq.~13 in the manuscript) uses absolute cosine similarity
  via representational dissimilarity; the $\Corr$ function is $C^\infty$
  only when the covariance matrix is full rank, which fails at degenerate
  activations (e.g., uniform hidden states across all inputs).
\end{enumerate}
The original proof's claim of $C^\infty$ smoothness is therefore incorrect.

\paragraph{Step 2: The vector field is locally Lipschitz.}
Despite non-smoothness at measure-zero sets, each component function is
piecewise-$C^\infty$ with bounded subgradients:
\begin{itemize}[nosep]
\item \textbf{Polynomial terms} ($\sigma_A \cdot \delta_A$, $\alpha_A(1-\sigma_A)$,
  etc.): globally Lipschitz on the compact set $\mathcal{X}$ with constant
  bounded by the maximum derivative of each monomial.
\item \textbf{Gompertz efficiency} $\eta(d,t) = \eta_{\max} \exp(-a \exp(-b \delta_A^{\text{relative}}))$:
  $C^\infty$ and bounded derivative on $\delta_A^{\text{relative}} \in [0,1]$.
\item \textbf{max}(0, $\cdot$): 1-Lipschitz with a single non-differentiable
  point at the origin.
\item \textbf{Pearson correlation} $\Corr(\cdot, \cdot)$: $C^\infty$ on
  the open set where the covariance matrix is positive definite; on the
  degenerate locus the function extends continuously with bounded
  difference quotients\footnote{Formally, if $\text{RDM}_{\text{rep}}$ is
  constant, the correlation is undefined. In practice, the Stage~1
  computation returns $r_A = 0$ in this case, making the function
  piecewise-$C^\infty$ with a jump discontinuity. The jump size is
  bounded by $O(1/\sqrt{n_{\text{batch}}})$ per Proposition~3.6.}.
\item \textbf{Cosine similarity} $\CosSim(\cdot,\cdot)$: $C^\infty$ away
  from zero vectors; 1-Lipschitz on the unit sphere.
\end{itemize}
The composition of Lipschitz functions is Lipschitz. Since $\mathcal{X}$
is compact, each component function $f_i$ has a finite Lipschitz constant
$L_i$ on $\text{int}(\mathcal{X})$.

\paragraph{Step 3: Apply Picard--Lindelöf.}
By Step~2, $\mathbf{f}$ is locally Lipschitz on $\text{int}(\mathcal{X})$.
The Picard--Lindelöf theorem guarantees a unique local solution
$\mathbf{x}_A(t)$ through any initial condition $\mathbf{x}_A(0) \in
\text{int}(\mathcal{X})$ on a maximal interval $[0, T_{\max})$ with
$T_{\max} > 0$. If $T_{\max} < \infty$, the solution must approach
$\partial\mathcal{X}$ as $t \to T_{\max}$ (the escape lemma).
$\square$
\end{proof}

% --- Boundary Cases ---
\begin{boundary-case}
\textbf{$\mathbf{x}_A(0) \in \partial\mathcal{X}$ (boundary initial
conditions):} Picard--Lindelöf requires open domain; on the boundary,
the vector field may not be Lipschitz (e.g., $\max(0,\cdot)$ at the
origin). Existence may still hold (e.g., via Carathéodory's theorem)
but uniqueness can fail. This is acceptable because the paper uses
$\mathbf{x}_A(0) \in \text{int}(\mathcal{X})$ for all practical
training scenarios.

\textbf{Zero determinant of $\Corr$ covariance:} If the RDM is
singular (e.g., all hidden states identical), the correlation is
undefined. The Stage~1 protocol must handle this by returning
$r_A = \text{undefined} \to 0$, introducing a single-point
discontinuity. Uniqueness at such points requires Filippov
solutions.
\end{boundary-case}

% --- Circular Dependency Check ---
\begin{circularity-check}
\textbf{Does Proposition~3.1 depend on any later result?} It requires
Assumption~\ref{assum:regularity} on boundedness and Lipschitz
continuity of coefficient functions. These are structural assumptions
about the training environment, not results derived later. 
\textbf{Proposition~3.6} provides the $O(1/\sqrt{n_{\text{batch}}})$
bound for the correlation estimation error, but it is not strictly
necessary for existence --- existence holds even without the bound.

\textbf{Does any later result depend on Proposition~3.1?}
Proposition~3.2 (forward invariance) requires the existence of a
unique solution to analyze its boundary behavior, creating a logical
dependency: 3.1 $\to$ 3.2. This is a forward dependency, not a
circular one. Safe.
\end{circularity-check}
